\documentclass[10pt]{article}

\usepackage{arxiv}
\usepackage[utf8]{inputenc}
\usepackage[T1]{fontenc}
\usepackage{hyperref}
\usepackage{url}
\usepackage{booktabs}
\usepackage{amsfonts}
\usepackage{amsmath}
\usepackage{graphicx}
\usepackage{enumitem}
\usepackage{listings}
\usepackage{xcolor}
\usepackage{longtable}
\usepackage{array}
\usepackage{multirow}
\usepackage{tabularx}
\usepackage{float}
\usepackage{tikz}

% Remove all colored frames from hyperlinks
\hypersetup{
  colorlinks=false,
  linkcolor=black,
  filecolor=black,
  citecolor=black,
  urlcolor=black,
  pdfborder={0 0 0}
}

% Code listing style
\lstdefinestyle{pythonstyle}{
  language=Python,
  basicstyle=\ttfamily\scriptsize,
  keywordstyle=\bfseries,
  commentstyle=\itshape,
  stringstyle=\ttfamily,
  frame=single,
  framerule=0.4pt,
  rulecolor=\color{black!30},
  backgroundcolor=\color{black!2},
  breaklines=true,
  breakatwhitespace=true,
  numbers=left,
  numberstyle=\tiny,
  numbersep=5pt,
  xleftmargin=15pt,
  showstringspaces=false,
  columns=flexible
}
\lstset{style=pythonstyle}

% Compact lists
\setlist{nosep, leftmargin=*}
\setlist[itemize]{label=\textbullet}
\setlist[enumerate]{label=\arabic*.}

\title{Federated Learning Physical AI Oncology Trials Unification}

\author{
  Kevin Kawchak \\
  CEO, ChemicalQDevice \\
  \href{https://orcid.org/0009-0007-5457-8667}{\textcolor{green!60!black}{\textbf{ORCID}}} \href{https://orcid.org/0009-0007-5457-8667}{https://orcid.org/0009-0007-5457-8667} \\
  \texttt{kevink@chemicalqdevice.com}
}

\date{February 27, 2026}

\begin{document}
\maketitle

\begin{abstract}
The transition from conventional software-only artificial intelligence to physical AI systems incorporating robotic hardware in oncology clinical trials represents a paradigm shift requiring unified infrastructure for privacy, regulation, cross-framework interoperability, and multi-organization cooperation. This paper presents the PAI Oncology Trial FL platform (v1.1.0), a comprehensive federated learning framework comprising 235 Python modules ($\sim$86,800 lines of code) that unifies five critical infrastructure pillars: (1) Privacy Infrastructure implementing all 18 HIPAA Safe Harbor identifiers with HMAC-SHA256 pseudonymization, (2) Regulatory Infrastructure spanning FDA, IRB, ICH-GCP, and multi-jurisdiction compliance across v0.6.0 and v0.9.1, (3) Cross-Framework Unification bridging NVIDIA Isaac Sim, MuJoCo, Gazebo, and PyBullet simulation environments, (4) Standards \& Benchmarking for Q1 2026 objectives including model conversion and registry pipelines, and (5) Multi-Organization Cooperation enabling federated training across academic medical centers, community hospitals, and pharmaceutical companies. End-to-end workflow demonstrations are presented across 31 example scripts, 6 agentic AI production examples implementing Model Context Protocol (MCP), ReAct reasoning, real-time monitoring, autonomous orchestration, safety-constrained execution, and RAG-based compliance. A triple AI peer review process (v0.9.4--v0.9.9) using sequential Codex-to-Claude Code review-fix cycles resolved 31/31 code recommendations at 100\% completion, establishing a dual-manufacturer trust benchmark for AI-generated clinical trial software. The platform demonstrates that unified federated learning infrastructure is a necessary precondition for transitioning the oncology industry to using robots in physical AI clinical trials.

\keywords{Federated Learning \and Physical AI \and Oncology Clinical Trials \and Digital Twins \and Regulatory Compliance \and HIPAA \and Cross-Framework Unification \and Agentic AI \and Model Context Protocol \and Differential Privacy}
\end{abstract}

\vspace{-0.5em}
\tableofcontents
\vspace{0.5em}

% =================================================================
\section{Introduction}
\label{sec:introduction}
% =================================================================

\subsection{Current AI in Federated Learning Oncology}

Federated learning (FL) has emerged as the dominant paradigm for privacy-preserving collaborative machine learning in oncology, enabling multiple institutions to jointly train AI models without exchanging raw patient data~\cite{mcmahan2017fedavg, rieke2020fl_medicine}. Leading cancer centers including Dana-Farber have launched federated AI platforms where models travel to each center's secure data environment~\cite{danafarberfederated2025}, and the European DigiONE consortium has demonstrated privacy-preserving learning health systems in precision oncology~\cite{digione2024}.

Current FL approaches in oncology typically address isolated aspects of the clinical trial workflow:
\begin{itemize}
  \item \textbf{Model training:} FedAvg~\cite{mcmahan2017fedavg}, FedProx~\cite{li2020fedprox}, and SCAFFOLD~\cite{karimireddy2020scaffold} for heterogeneous hospital data
  \item \textbf{Privacy:} Differential privacy with Gaussian noise mechanisms~\cite{dwork2014dp} and secure aggregation protocols
  \item \textbf{Brain tumor segmentation:} Multi-institutional FL for glioblastoma without patient data sharing~\cite{sheller2020fl_brain}
\end{itemize}

However, no existing open-source platform unifies federated learning, physical AI robotics, digital twins, regulatory compliance, privacy infrastructure, and clinical analytics into a single coherent framework for oncology clinical trials~\cite{kawchak2026fl}.

\subsection{Benefits of Upgrading to AI Robots in Oncology Trials}

The integration of physical AI---embodied robotic systems with onboard intelligence---into oncology clinical trials offers transformative benefits:
\begin{enumerate}
  \item \textbf{Precision in tumor biopsies and drug delivery:} Surgical robots executing procedures with sub-millimeter accuracy guided by digital twin predictions, as modeled in the PAI platform's \texttt{physical\_ai/} modules (v0.1.0)
  \item \textbf{Consistent monitoring:} Autonomous robotic systems performing 24/7 patient monitoring with real-time anomaly detection via multi-modal sensor fusion (v0.3.0)
  \item \textbf{Cross-site standardization:} Unified robot programming across trial sites eliminates procedural variability, enabled by the cross-framework unification layer (v0.3.0)
  \item \textbf{Federated model improvement:} Robotic telemetry from multiple sites feeds back into federated model training without raw data exchange, as demonstrated in \texttt{01\_mcp\_oncology\_server.py} (v0.5.0)
\end{enumerate}

\subsection{Projected Physical AI Benefits from Quantitative Codebase Data}

The PAI Oncology Trial FL platform (v1.1.0) provides quantitative evidence for physical AI readiness:

\begin{table}[H]
\centering
\caption{Platform quantitative metrics supporting physical AI oncology trial readiness}
\label{tab:quantitative}
\small
\begin{tabular}{@{}lrl@{}}
\toprule
\textbf{Metric} & \textbf{Value} & \textbf{Version} \\
\midrule
Total Python modules & 235 & v1.0.1 \\
Lines of code & $\sim$86,800 & v1.0.1 \\
Test files & 82 & v0.8.0+ \\
Test directories & 15 & v0.8.0+ \\
Example scripts & 31 & v0.5.0 \\
Agentic AI examples & 6 & v0.5.0 \\
CLI tools & 5 & v0.4.0 \\
Security audit findings resolved & 61 & v0.7.0 \\
Peer review recommendations resolved & 31/31 (100\%) & v0.9.4--v0.9.9 \\
Regulatory frameworks supported & 9 & v0.9.1 \\
HIPAA Safe Harbor identifiers & 18/18 & v0.6.0 \\
Simulation frameworks bridged & 4 & v0.3.0 \\
Tumor growth models & 4 & v0.4.0 \\
Supported Python versions & 3.10, 3.11, 3.12 & v0.3.0+ \\
\bottomrule
\end{tabular}
\end{table}

The upstream Physical AI Oncology Trials repository~\cite{kawchak2026physical} introduced the Unification Standard Level (USL) scoring framework evaluating 9 robotic systems across humanoid robots, surgical robots, and collaborative robots, with scores ranging from 3.4 (UFACTORY xArm7) to 7.4 (Franka Panda)~\cite{kawchak2026usl}. The PAI Oncology Trial FL platform extends this foundation with federated learning capabilities necessary for multi-site robot deployment.

% =================================================================
\section{Methods}
\label{sec:methods}
% =================================================================

\subsection{AI Development Tools and Models}

The PAI Oncology Trial FL platform was developed using the following AI systems:
\begin{itemize}
  \item \textbf{Primary development:} Anthropic Claude Code powered by Claude Opus 4.6~\cite{anthropic2026claude}---responsible for all code generation, documentation, and fix implementation across v0.1.0 through v1.1.0
  \item \textbf{Independent code review:} OpenAI ChatGPT-5.2-Codex~\cite{openai2026codex}---performed three sequential senior engineering review cycles (v0.9.4, v0.9.6, v0.9.8) generating 31 total recommendations
  \item \textbf{Paper generation:} This paper was primarily generated by Claude Code (Opus 4.6) based on Mr.\ Kawchak's hand-written prompt, which took approximately 2 hours to compose
\end{itemize}

\subsection{Development Timeline and Version History}

The codebase was developed over the period February 17--27, 2026, with the repository converted from private to public on February 27, 2026. Table~\ref{tab:timeline} summarizes the complete version history.

\begin{table}[H]
\centering
\caption{Complete release timeline with AI contributor roles}
\label{tab:timeline}
\small
\begin{tabular}{@{}llll@{}}
\toprule
\textbf{Version} & \textbf{Date} & \textbf{AI Role} & \textbf{Key Deliverable} \\
\midrule
v0.1.0 & 2026-02-17 & Claude Code & Initial FL platform \\
v0.2.0 & 2026-02-17 & Claude Code & FedProx, SCAFFOLD, data harmonization \\
v0.3.0 & 2026-02-17 & Claude Code & Unification framework, CI/CD \\
v0.4.0 & 2026-02-17 & Claude Code & Digital twins, 5 CLI tools \\
v0.5.0 & 2026-02-17 & Claude Code & 31 example scripts \\
v0.6.0 & 2026-02-18 & Claude Code & Privacy \& regulatory (9 modules) \\
v0.7.0 & 2026-02-18 & Claude Code & Security audit (61 findings) \\
v0.8.0 & 2026-02-18 & Claude Code & Test suite (82 files, 15 dirs) \\
v0.9.0 & 2026-02-18 & Claude Code & Clinical analytics (7 modules) \\
v0.9.1 & 2026-02-18 & Claude Code & Regulatory submissions (6 modules) \\
v0.9.2 & 2026-02-18 & Claude Code & Release documentation \\
v0.9.3 & 2026-02-19 & Claude Code & Visualization (30 charts) \\
v0.9.4 & 2026-02-19 & Codex & Peer review \#1: 12 recommendations \\
v0.9.5 & 2026-02-19 & Claude Code & Fix cycle \#1: 12/12 resolved \\
v0.9.6 & 2026-02-19 & Codex & Peer review \#2: 9 recommendations \\
v0.9.7 & 2026-02-19 & Claude Code & Fix cycle \#2: 9/9 resolved \\
v0.9.8 & 2026-02-19 & Codex & Peer review \#3: 10 recommendations \\
v0.9.9 & 2026-02-19 & Claude Code & Fix cycle \#3: 10/10 resolved \\
v1.0.0 & 2026-02-19 & Claude Code & Stable release \\
v1.0.1 & 2026-02-19 & Claude Code & Documentation update \\
v1.1.0 & 2026-02-27 & Claude Code & This paper and archival \\
\bottomrule
\end{tabular}
\end{table}

\textbf{Development statistics:}
\begin{itemize}
  \item \textbf{Total span:} 11 days (Feb 17--27, 2026)
  \item \textbf{Days with AI code updates:} 4 days (Feb 17, 18, 19, 27)
  \item \textbf{Days of AI code recommendations (Codex):} 1 day (Feb 19---v0.9.4, v0.9.6, v0.9.8)
  \item \textbf{Days of AI code fixes (Claude Code):} 4 days (Feb 17, 18, 19, 27---v0.1.0 through v1.1.0)
  \item \textbf{Human prompt authoring (Kevin Kawchak):} 18 prompts across development lifecycle, including a 2-hour hand-written prompt for this paper
  \item \textbf{Time from first prompt to PR submission:} Approximately 30 minutes (prompt submission: 2026-02-27 03:15 UTC, PR creation: 2026-02-27 03:45 UTC est.)
\end{itemize}

\subsection{Prompt-Driven Development Methodology}

All 18 development prompts are recorded in \texttt{prompts.md}~\cite{kawchak2026fl}. The methodology follows a single-prompt delivery pattern where each prompt generates a complete, self-contained release. Key prompts include:
\begin{itemize}
  \item \textbf{Prompt 1 (v0.1.0):} ``Integrate physical AI with federated learning platform instructions into \texttt{pai-oncology-trial-fl}. Utilize \texttt{physical-ai-oncology-trials} for physical AI implementation.''
  \item \textbf{Prompt 8 (v0.8.0):} ``Build comprehensive pytest-based test suite achieving full coverage of all Python modules.''
  \item \textbf{Prompt 13a (v0.9.4):} ``Provide comprehensive code reviews to find errors across the entire codebase'' (Codex)
  \item \textbf{Prompt 16 (v1.0.0):} ``Prove and push a v1.0.0 stable release by analyzing the codebase end-to-end''
\end{itemize}

% =================================================================
\section{Results}
\label{sec:results}
% =================================================================

\subsection{Platform Architecture Overview}

The PAI Oncology Trial FL platform follows a layered architecture where five infrastructure pillars support the end-to-end clinical trial workflow. Table~\ref{tab:architecture} maps each pillar to its codebase location, version of introduction, and quantitative scope.

\begin{table}[H]
\centering
\caption{Five-pillar architecture with codebase metrics}
\label{tab:architecture}
\small
\begin{tabular}{@{}clcrl@{}}
\toprule
\textbf{\#} & \textbf{Pillar} & \textbf{Directory} & \textbf{LOC} & \textbf{Version} \\
\midrule
1 & Privacy Infrastructure & \texttt{privacy/} & $\sim$5,400 & v0.6.0 \\
2 & Regulatory Infrastructure & \texttt{regulatory/}, \texttt{reg-sub/} & $\sim$10,000 & v0.6.0, v0.9.1 \\
3 & Cross-Framework Unification & \texttt{unification/} & $\sim$8,200 & v0.3.0 \\
4 & Standards \& Benchmarking & \texttt{q1-2026-standards/} & $\sim$3,800 & v0.3.0 \\
5 & Multi-Org Cooperation & \texttt{federated/} & $\sim$12,000 & v0.1.0--v0.2.0 \\
\midrule
  & Supporting: Clinical Analytics & \texttt{clinical-analytics/} & $\sim$4,900 & v0.9.0 \\
  & Supporting: Digital Twins & \texttt{digital-twins/} & $\sim$4,500 & v0.4.0 \\
  & Supporting: Physical AI & \texttt{physical\_ai/} & $\sim$6,200 & v0.1.0 \\
  & Supporting: Agentic AI & \texttt{agentic-ai/} & $\sim$6,700 & v0.5.0 \\
  & Infrastructure: Tests & \texttt{tests/} & $\sim$15,000 & v0.8.0 \\
  & Infrastructure: CLI Tools & \texttt{tools/} & $\sim$5,100 & v0.4.0 \\
  & Infrastructure: Utilities & \texttt{utils/}, \texttt{scripts/} & $\sim$500 & v0.9.7--v0.9.9 \\
\bottomrule
\end{tabular}
\end{table}

\begin{verbatim}
  PAI Oncology Trial FL - End-to-End Architecture (v1.1.0)
  ========================================================

  +-- Multi-Organization Cooperation (Pillar 5) ----------+
  |  FedAvg | FedProx | SCAFFOLD | Secure Aggregation     |
  |  Site Enrollment | Data Harmonization | DP Budget      |
  +-------+------+------+------+------+------+------------+
          |      |      |      |      |
  +-------+--+ +-+------+-+ +--+------+--+ +-+----------+
  | Privacy  | | Regulatory | | Unification | | Standards |
  | (Pillar  | | (Pillar 2) | | (Pillar 3)  | | (Pillar  |
  |  1)      | |            | |             | |  4)       |
  | PHI Det. | | FDA Track  | | Isaac-MuJo  | | Model     |
  | De-ID    | | IRB Mgmt   | | Bridge      | | Registry  |
  | RBAC     | | GCP Check  | | Agent Unif. | | Benchmark |
  | Breach   | | eCTD Comp. | | Format Conv | | Convert   |
  +----+-----+ +-----+------+ +------+------+ +-----+----+
       |             |               |               |
  +----+-------------+---------------+---------------+----+
  |        Supporting Platform Modules                     |
  | Digital Twins | Clinical Analytics | Agentic AI | CLI  |
  | PK/PD | Survival | Risk | DSMB | MCP | ReAct | RAG   |
  +-------------------------------------------------------+
          |
  +-------+-----------------------------------------------+
  |              Tests & Quality Infrastructure            |
  | 82 test files | 15 directories | Security Audit       |
  | Triple AI Peer Review | CI (3.10/3.11/3.12)           |
  +-------------------------------------------------------+
\end{verbatim}

\subsection{Pillar 1: Privacy Infrastructure (v0.6.0)}
\label{sec:privacy}

The privacy framework implements comprehensive HIPAA Safe Harbor compliance across five modules totaling $\sim$5,400 LOC within the \texttt{privacy/} directory.

\subsubsection{PHI/PII Detection and Management}

The \texttt{phi\_detector.py} module (v0.6.0) implements all 18 HIPAA Safe Harbor identifiers per 45 CFR 164.514(b)(2)~\cite{hipaa_safeharbor}:

\begin{lstlisting}[caption={PHI detection categories from \texttt{privacy/phi-pii-management/phi\_detector.py} (v0.6.0)}]
class PHICategory(Enum):
    NAME = "name"
    GEOGRAPHIC = "geographic_subdivision"
    DATE = "dates"
    PHONE = "phone_number"
    FAX = "fax_number"
    EMAIL = "email_address"
    SSN = "social_security_number"
    MEDICAL_RECORD = "medical_record_number"
    HEALTH_PLAN = "health_plan_beneficiary"
    ACCOUNT = "account_number"
    LICENSE = "certificate_license"
    VEHICLE = "vehicle_identifier"
    DEVICE = "device_identifier"
    URL = "web_url"
    IP_ADDRESS = "ip_address"
    BIOMETRIC = "biometric_identifier"
    PHOTO = "full_face_photo"
    OTHER = "other_unique_identifier"
\end{lstlisting}

\subsubsection{De-identification Pipeline}

The de-identification module provides six transformation methods with HMAC-SHA256 pseudonymization using cryptographically random salts (\texttt{os.urandom(32)}), as hardened during peer review cycle 3 (v0.9.8/v0.9.9):

\begin{itemize}
  \item \textbf{REDACT:} Complete value removal
  \item \textbf{HASH:} HMAC-SHA256 pseudonymization (environment-driven keys post-v0.9.9)
  \item \textbf{GENERALIZE:} Geographic data to 3-digit ZIP codes
  \item \textbf{SUPPRESS:} Selective field suppression
  \item \textbf{PSEUDONYMIZE:} Consistent patient-level token generation
  \item \textbf{DATE\_SHIFT:} Consistent patient-level date offsets
\end{itemize}

\subsubsection{Access Control and Audit}

Role-based access control (RBAC) implements 21 CFR Part 11 compliant audit trails~\cite{cfr_part11} with five role tiers, time-limited access with fail-closed expiration (v0.6.0), and immutable audit log copies returned by \texttt{get\_audit\_log()} (v0.7.0 security fix).

\subsubsection{Breach Response and Data Use Agreements}

The breach response protocol implements four-factor risk assessment per 45 CFR 164.402 with 60-day notification tracking. The DUA generator supports five agreement types (research, clinical, vendor, multi-institutional, pediatric) with security tiers and configurable retention policies.

\subsection{Pillar 2: Regulatory Infrastructure (v0.6.0, v0.9.1)}
\label{sec:regulatory}

\subsubsection{Core Regulatory Framework (v0.6.0)}

The \texttt{regulatory/} directory implements five compliance domains:

\begin{table}[H]
\centering
\caption{Regulatory infrastructure modules and compliance alignment}
\label{tab:regulatory}
\small
\begin{tabular}{@{}lll@{}}
\toprule
\textbf{Module} & \textbf{Version} & \textbf{Standards} \\
\midrule
\texttt{fda\_submission\_tracker.py} & v0.6.0 & FDA 510(k), De Novo, PMA \\
\texttt{irb\_protocol\_manager.py} & v0.6.0 & ICH E6(R3), 21 CFR Part 11 \\
\texttt{gcp\_compliance\_checker.py} & v0.6.0 & ICH E6(R3) 13 principles \\
\texttt{regulatory\_tracker.py} & v0.6.0 & FDA, EMA, PMDA, TGA, HC \\
\texttt{HUMAN\_OVERSIGHT\_QMS.md} & v0.6.0 & CRF risk tiers, AE boundaries \\
\bottomrule
\end{tabular}
\end{table}

The FDA submission tracker supports five regulatory pathways---510(k), De Novo, PMA, Breakthrough, and Pre-Submission---referencing FDA AI/ML Device Guidance (January 2025)~\cite{fda2025aiml}, PCCP Guidance (August 2025)~\cite{fda2025pccp}, and QMSR (February 2026). The GCP compliance checker scores 13 ICH E6(R3) principles~\cite{ich2025e6r3} excluding \texttt{NOT\_ASSESSED} findings from the denominator (v0.7.0 logic fix).

\subsubsection{Regulatory Submissions Platform (v0.9.1)}

The \texttt{regulatory-submissions/} platform ($\sim$4,600 LOC across 6 modules) complements clinical analytics with downstream regulatory packaging:

\begin{enumerate}
  \item \textbf{Submission Orchestrator:} Lifecycle management for FDA submission workflows
  \item \textbf{eCTD Compiler:} Module structure generation per FDA Technical Conformance Guide with SHA-256 checksums
  \item \textbf{Compliance Validator:} Multi-regulation validation across 9 frameworks (HIPAA, GDPR, FDA 21 CFR 820, Part 11, ISO 14971, IEC 62304, ICH E6(R3)/E9(R1), MDR 2017/745)
  \item \textbf{Document Generator:} Automated CSR, SAP, PCCP, ISO 14971, IEC 62304 documents
  \item \textbf{Regulatory Intelligence:} 7-jurisdiction guidance tracking (FDA, EMA, PMDA, TGA, Health Canada, MHRA, NMPA)
  \item \textbf{Submission Analytics:} KPIs, trend analysis, FDA MDUFA benchmarking
\end{enumerate}

\subsection{Pillar 3: Cross-Framework Unification (v0.3.0)}
\label{sec:unification}

The \texttt{unification/} directory (v0.3.0) enables seamless interoperability across four simulation frameworks for surgical robotics:

\begin{verbatim}
  Cross-Framework Unification Architecture (v0.3.0)
  ==================================================

  +------------------+       +------------------+
  | NVIDIA Isaac Sim |<----->|     MuJoCo       |
  | (High-fidelity)  |       |  (Fast physics)  |
  +--------+---------+       +--------+---------+
           |    Bidirectional Bridge   |
           +----------+  +------------+
                      |  |
               +------+--+------+
               | Unified State   |
               | Representation  |
               +------+--+------+
                      |  |
           +----------+  +------------+
           |                          |
  +--------+---------+       +--------+---------+
  |     Gazebo       |       |    PyBullet      |
  |   (ROS 2)        |       | (Lightweight)    |
  +------------------+       +------------------+
\end{verbatim}

Key unification components include:
\begin{itemize}
  \item \textbf{Isaac-MuJoCo Bridge} (\texttt{simulation\_physics/}): Bidirectional state conversion with parameter mapping, enabling policy transfer between high-fidelity and fast physics engines
  \item \textbf{Unified Agent Interface} (\texttt{agentic\_generative\_ai/}): Tool format conversion for CrewAI, LangGraph, AutoGen, and custom backends
  \item \textbf{Cross-Platform Tools} (\texttt{cross\_platform\_tools/}): Framework detection (Isaac Lab, MuJoCo, Gazebo, PyBullet), model format conversion (URDF, MJCF, SDF, USD), and policy export with validation
  \item \textbf{Standards Protocols} (\texttt{standards\_protocols/}): Data format standards, communication protocols, and safety compliance
\end{itemize}

\subsection{Pillar 4: Standards \& Benchmarking (v0.3.0)}
\label{sec:standards}

The \texttt{q1-2026-standards/} directory implements three Q1 2026 objectives:

\begin{enumerate}
  \item \textbf{Objective 1---Model Conversion Pipeline:} Cross-framework model format conversion with validation, supporting URDF$\leftrightarrow$MJCF$\leftrightarrow$SDF$\leftrightarrow$USD bidirectional transforms
  \item \textbf{Objective 2---Model Registry:} Federated model registry with validation gates, version tracking, and compliance metadata
  \item \textbf{Objective 3---Cross-Platform Benchmarking:} Standardized benchmark runner enabling reproducible comparison across simulation frameworks
\end{enumerate}

The implementation guide includes a timeline (\texttt{timeline.md}) and compliance checklist (\texttt{compliance\_checklist.md}) aligned with IEC 80601-2-77~\cite{iec80601} and ISO 14971~\cite{iso14971}.

\subsection{Pillar 5: Multi-Organization Cooperation (v0.1.0--v1.0.1)}
\label{sec:cooperation}

The federated architecture enables multi-organization cooperation across the oncology trial ecosystem:

\begin{table}[H]
\centering
\caption{Multi-organization cooperation model with platform integration points}
\label{tab:cooperation}
\small
\begin{tabular}{@{}p{3.2cm}p{3.8cm}l@{}}
\toprule
\textbf{Organization} & \textbf{Integration Points} & \textbf{Version} \\
\midrule
Academic Medical Centers & Federated training, IRB, digital twins & v0.1.0+ \\
Community Hospitals & Site enrollment, data harmonization & v0.2.0+ \\
Pharmaceutical Cos. & Treatment simulation, regulatory & v0.6.0+ \\
Robotics Manufacturers & Simulation frameworks, surgical tasks & v0.3.0+ \\
Regulatory Bodies & Submission tracking, compliance & v0.9.1+ \\
Standards Organizations & IEC/ISO/ICH alignment & v0.3.0+ \\
\bottomrule
\end{tabular}
\end{table}

The federated learning core (\texttt{federated/}) implements three aggregation strategies with privacy guarantees:
\begin{itemize}
  \item \textbf{FedAvg} (v0.1.0)~\cite{mcmahan2017fedavg}: Standard federated averaging for IID data distributions
  \item \textbf{FedProx} (v0.2.0)~\cite{li2020fedprox}: Proximal term regularization for non-IID heterogeneous hospital data
  \item \textbf{SCAFFOLD} (v0.2.0)~\cite{karimireddy2020scaffold}: Control variate correction for client drift across sites
  \item \textbf{Differential Privacy} (v0.1.0)~\cite{dwork2014dp}: Gaussian mechanism with gradient clipping and privacy budget ($\varepsilon$, $\delta$) tracking
  \item \textbf{Secure Aggregation} (v0.1.0): Mask-based protocol ensuring no raw model weights are exchanged
\end{itemize}

\subsection{Workflow Demonstrations: Domain Examples (v0.5.0)}
\label{sec:examples}

Version 0.5.0 introduced 31 production-quality example scripts across five domains, demonstrating end-to-end workflow capabilities. These examples serve as reference implementations that clinicians, engineers, and regulators can use to evaluate the platform.

\subsubsection{Core Federated Learning Examples}

Five core examples (\texttt{examples/}, 4,924 LOC total) demonstrate progressive complexity:
\begin{enumerate}
  \item \texttt{01\_federated\_training\_workflow.py} (881 LOC): Multi-site FedAvg/FedProx pipeline with DP integration, secure aggregation, and convergence monitoring
  \item \texttt{02\_digital\_twin\_planning.py} (924 LOC): Patient digital twin construction with exponential, logistic, Gompertz, and von Bertalanffy tumor growth models
  \item \texttt{03\_cross\_framework\_validation.py} (865 LOC): Policy validation across Isaac Sim, MuJoCo, PyBullet, and Gazebo with kinematic drift analysis
  \item \texttt{04\_agentic\_clinical\_workflow.py} (917 LOC): Multi-agent clinical decision support with oncologist, radiologist, safety, and data agents
  \item \texttt{05\_outcome\_prediction\_pipeline.py} (982 LOC): Federated outcome prediction with survival analysis, biomarker integration, and fairness auditing
\end{enumerate}

\subsubsection{Agentic AI Examples---Production Architecture}

Six agentic AI examples (\texttt{agentic-ai/examples-agentic-ai/}, $\sim$6,700 LOC total) provide production-grade implementations:

\begin{table}[H]
\centering
\caption{Agentic AI example scripts with architecture patterns (v0.5.0)}
\label{tab:agentic}
\small
\begin{tabular}{@{}clp{5cm}@{}}
\toprule
\textbf{\#} & \textbf{Script} & \textbf{Architecture Pattern} \\
\midrule
01 & \texttt{mcp\_oncology\_server.py} & Model Context Protocol server with tool registration, patient lookup, treatment simulation, protocol compliance, robotic telemetry \\
02 & \texttt{react\_treatment\_planner.py} & ReAct reasoning loops with thought-action-observation cycles, 10+ step traces \\
03 & \texttt{realtime\_adaptive\_monitoring.py} & Streaming multi-modal data, adaptive thresholds, CTCAE v5.0 grading \\
04 & \texttt{autonomous\_simulation\_orchestrator.py} & Multi-framework job coordination with Pareto optimization \\
05 & \texttt{safety\_constrained\_executor.py} & IEC 80601-2-77 safety gates, human-in-the-loop, emergency stop \\
06 & \texttt{oncology\_rag\_compliance.py} & RAG with semantic chunking, cross-encoder reranking, citation tracking \\
\bottomrule
\end{tabular}
\end{table}

\textbf{MCP Oncology Server (\texttt{01\_mcp\_oncology\_server.py}).} This script implements a Model Context Protocol server~\cite{mcp2025} exposing five oncology domain tools for LLM agents:

\begin{lstlisting}[caption={MCP tool registration pattern from \texttt{01\_mcp\_oncology\_server.py} (v0.5.0)}]
# Tool definitions with JSON Schema input validation
TOOLS = [
    {"name": "patient_lookup",
     "description": "Retrieve de-identified patient record",
     "inputSchema": {"type": "object",
       "properties": {"patient_id": {"type": "string"}}}},
    {"name": "treatment_simulation",
     "description": "Run digital twin treatment simulation",
     "inputSchema": {"type": "object",
       "properties": {"patient_id": {"type": "string"},
                      "modality": {"type": "string"}}}},
    {"name": "protocol_compliance",
     "description": "Verify ICH E6(R3) compliance",
     "inputSchema": {"type": "object",
       "properties": {"protocol_id": {"type": "string"}}}},
]
\end{lstlisting}

The MCP server integrates with the broader platform by:
\begin{itemize}
  \item Querying the \texttt{privacy/} de-identification pipeline before returning patient records
  \item Invoking the \texttt{digital-twins/} patient modeling engine for treatment simulation
  \item Validating against the \texttt{regulatory/} GCP compliance checker for protocol verification
  \item Logging all tool invocations to 21 CFR Part 11 compliant audit trails with hash-chain integrity
  \item Ingesting robotic procedure telemetry from physical AI surgical systems
\end{itemize}

\textbf{Synergistic inter-script operation.} The six agentic AI scripts function both independently and synergistically:
\begin{itemize}
  \item Scripts 01 (MCP) + 05 (Safety) combine to provide safety-constrained tool execution for LLM agents operating robotic systems
  \item Scripts 02 (ReAct) + 06 (RAG) combine reasoning with regulatory evidence grounding
  \item Scripts 03 (Monitoring) + 04 (Orchestration) enable real-time adaptive simulation campaigns where monitoring alerts trigger simulation parameter adjustments
  \item All six scripts share the platform's audit logging infrastructure, differential privacy protections, and HIPAA-compliant data handling
\end{itemize}

\subsubsection{Physical AI Examples}

Six physical AI examples (\texttt{examples-physical-ai/}, 800+ LOC each) demonstrate robotic oncology capabilities: real-time safety monitoring with ISO 13482 alignment, multi-sensor Kalman fusion, ROS 2 surgical deployment with lifecycle nodes, hand-eye calibration via Tsai-Lenz, shared autonomy teleoperation with variable blending, and robotic sample handling with chain-of-custody tracking.

\subsection{Triple AI Peer Review Pipeline (v0.9.4--v0.9.9)}
\label{sec:peerreview}

The triple AI peer review process constitutes a \emph{sequential} pipeline of review-fix cycles alternating between two AI manufacturers, establishing a dual-manufacturer trust benchmark for AI-generated code~\cite{kawchak2026peerreview}:

\begin{verbatim}
 Sequential Triple AI Peer Review (v0.9.4 - v0.9.9)
 ====================================================

 Claude Code --> Codex ------> Claude Code --> Codex
 (codebase)     (review #1)    (fixes #1)     (review #2)
  v0.1-0.9.3     v0.9.4         v0.9.5         v0.9.6
                  12 recs        12/12          9 recs

 --> Claude Code --> Codex ------> Claude Code
     (fixes #2)     (review #3)    (fixes #3)
      v0.9.7         v0.9.8         v0.9.9
      9/9            10 recs        10/10

 Total: 31/31 recommendations resolved (100%)
\end{verbatim}

This was \textbf{not} a parallel process. Each cycle was sequential: Claude Code provided the codebase, Codex reviewed and generated recommendations, Claude Code implemented all fixes, then Codex reviewed the updated codebase again. This sequential alternation---Claude Code $\to$ Codex $\to$ Claude Code $\to$ Codex $\to$ Claude Code $\to$ Codex---ensures each review cycle builds on the fixes from the previous cycle, enabling progressively deeper risk detection.

\begin{table}[H]
\centering
\caption{Triple AI peer review cycle metrics (v0.9.4--v0.9.9)}
\label{tab:peerreview}
\small
\begin{tabular}{@{}cccccc@{}}
\toprule
\textbf{Cycle} & \textbf{Review} & \textbf{Fixes} & \textbf{Recs} & \textbf{Resolved} & \textbf{Focus} \\
\midrule
1 & v0.9.4 (Codex) & v0.9.5 (Claude) & 12 & 12/12 & CI \& process \\
2 & v0.9.6 (Codex) & v0.9.7 (Claude) & 9 & 9/9 & Compliance \& security \\
3 & v0.9.8 (Codex) & v0.9.9 (Claude) & 10 & 10/10 & Hardening \& secrets \\
\midrule
\textbf{Total} & \textbf{3 reviews} & \textbf{3 fix cycles} & \textbf{31} & \textbf{31/31} & \textbf{100\%} \\
\bottomrule
\end{tabular}
\end{table}

\subsubsection{v1.0.1: Claude Code + Codex---31/31 Resolved as AI Bias Mitigation}

The 31/31 resolution rate using two different AI manufacturers (Anthropic for code generation/fixes, OpenAI for independent review) serves as an example of AI bias mitigation in software development. By ensuring neither manufacturer's blind spots persist in the final codebase, the dual-manufacturer approach establishes an emerging software quality methodology:
\begin{itemize}
  \item \textbf{Manufacturer A (Anthropic Claude Code):} May have systematic blind spots in security patterns, exception handling, or compliance areas
  \item \textbf{Manufacturer B (OpenAI Codex):} Independent review catches issues that Manufacturer A's training data may not flag
  \item \textbf{Result:} The alternating sequential process eliminates single-manufacturer bias, producing code that has been independently validated by two separate AI architectures
\end{itemize}

\subsection{Code Trust and Safety: Tests, Checks, and Migrations}
\label{sec:codetrust}

\subsubsection{Core Test Suite (v0.8.0)}

The test suite at \texttt{tests/} (v0.8.0) contains 82 test files across 15 directories, constituting a core infrastructure for code trust:

\begin{table}[H]
\centering
\caption{Test suite directory structure (v0.8.0, 15 directories)}
\label{tab:tests}
\small
\begin{tabular}{@{}lcp{4.5cm}@{}}
\toprule
\textbf{Directory} & \textbf{Files} & \textbf{Coverage Scope} \\
\midrule
\texttt{test\_federated/} & 6 & Coordinator, client, aggregation, DP \\
\texttt{test\_privacy/} & 6 & PHI, de-identification, RBAC, breach \\
\texttt{test\_regulatory/} & 5 & FDA, IRB, GCP, intelligence, oversight \\
\texttt{test\_clinical\_analytics/} & 7 & Orchestration, PK/PD, survival, DSMB \\
\texttt{test\_regulatory\_submissions/} & 4 & eCTD, compliance, documents, analytics \\
\texttt{test\_agentic\_ai/} & 5 & MCP, ReAct, monitoring, orchestrator \\
\texttt{test\_digital\_twins/} & 3 & Patient modeling, treatment, integration \\
\texttt{test\_physical\_ai/} & 4 & Digital twin, robotics, sensor fusion \\
\texttt{test\_unification/} & 3 & Bridge, agent interface, converter \\
\texttt{test\_tools/} & 2 & CLI tool validation \\
\texttt{test\_standards/} & 3 & Q1 2026 conversion, registry, benchmark \\
\texttt{test\_examples/} & 3 & Example script execution \\
\texttt{test\_regression/} & 2 & v0.7.0 audit finding regression guards \\
\texttt{test\_integration/} & 2 & End-to-end cross-module pipelines \\
\texttt{conftest.py} & 1 & Shared fixtures, \texttt{np.random.seed(42)} \\
\midrule
\textbf{Total} & \textbf{82} (incl.\ 16 \texttt{\_\_init\_\_.py}) & \\
\bottomrule
\end{tabular}
\end{table}

\subsubsection{Security Audit (v0.7.0)}

The comprehensive security audit (v0.7.0) identified and resolved 61 findings across three categories:
\begin{itemize}
  \item \textbf{12 Security vulnerabilities:} \texttt{torch.load} $\to$ \texttt{weights\_only=True}, \texttt{np.load} $\to$ \texttt{allow\_pickle=False}, plaintext key removal, mutable audit log exposure
  \item \textbf{14 Logic bugs:} Division-by-zero guards, bounded \texttt{while} loops with \texttt{MAX\_ITERATIONS}, floating-point tolerance (\texttt{np.isclose}), hazard ratio sign correction
  \item \textbf{35 Compliance findings:} \texttt{RESEARCH USE ONLY} disclaimers in all modules, audit trail logging, date shift handling, \texttt{model\_type="unspecified"} defaults
\end{itemize}

\subsubsection{Peer Review Fixes and Migrations}

\textbf{Cycle 1 (v0.9.4$\to$v0.9.5)---CI Determinism:}
\begin{itemize}
  \item Created \texttt{constraints/lint.txt} with pinned tool versions for CI reproducibility
  \item Added \texttt{.pre-commit-config.yaml} for local lint consistency
  \item 5 files modified, 4 new files, +67 net LOC
\end{itemize}

\textbf{Cycle 2 (v0.9.6$\to$v0.9.7)---Compliance \& Security:}
\begin{itemize}
  \item Fixed 24 UTC timestamp instances across 23 files (migrated to \texttt{datetime.timezone.utc})
  \item Narrowed 17 broad exception handlers to specific exception types
  \item Replaced MD5 hashing with SHA-256
  \item Created \texttt{scripts/check\_version\_alignment.py} (71 LOC), \texttt{scripts/security\_scan.py} (96 LOC), \texttt{scripts/release\_metrics.py} (89 LOC)
  \item 23 files modified, 3 new files, +25 net LOC
\end{itemize}

\textbf{Cycle 3 (v0.9.8$\to$v0.9.9)---Hardening \& Secrets:}
\begin{itemize}
  \item Removed hardcoded cryptographic keys from 8 modules (migrated to \texttt{os.environ} via \texttt{utils/crypto.py})
  \item Implemented audit integrity chain verification
  \item Created PHI-safe logging sanitizer (\texttt{utils/log\_sanitizer.py}, 36 LOC)
  \item Narrowed 11 additional exception handlers
  \item Created centralized configuration (\texttt{utils/config.py}, 66 LOC) and structured error codes (\texttt{utils/error\_codes.py}, 74 LOC)
  \item Added PHI fuzz tests (\texttt{tests/test\_privacy/test\_phi\_fuzz.py}, 127 LOC)
  \item 22 files modified, 7 new files, +386 net LOC
\end{itemize}

\textbf{Cumulative impact across all trust \& safety activities:}
\begin{itemize}
  \item 92 total issues resolved (61 audit + 31 peer review)
  \item 28 exception handlers narrowed across 19 files
  \item 25 UTC timestamp fixes
  \item 13 new utility scripts/modules
  \item +478 net lines of trust infrastructure code
\end{itemize}

\subsection{Clinical Analytics Platform (v0.9.0)}

The \texttt{clinical-analytics/} platform ($\sim$4,900 LOC across 7 modules) provides privacy-preserving analytics for multi-site oncology trials:

\begin{itemize}
  \item \textbf{PK/PD Engine:} One- and two-compartment pharmacokinetic models with analytical and ODE solutions (\texttt{scipy.integrate.solve\_ivp}), NCA metrics (AUC, C$_{\text{max}}$, T$_{\text{max}}$, t$_{1/2}$), E$_{\text{max}}$ dose-response
  \item \textbf{Survival Analysis:} Kaplan-Meier with Greenwood SE, log-rank test, Cox proportional hazards via BFGS with Breslow ties, Harrell's C-index, restricted mean survival time (RMST)
  \item \textbf{Risk Stratification:} Composite weighted scores, adaptive enrichment with per-stratum decision support, Hosmer-Lemeshow calibration
  \item \textbf{Clinical Interoperability:} ICD-10/SNOMED CT crosswalk, LOINC, RxNorm, MedDRA mappings, CDISC SDTM export~\cite{cdisc2024sdtm}
  \item \textbf{Consortium Reporting:} DSMB reports with SHA-256 integrity hashing for tamper detection
\end{itemize}

\subsection{Digital Twins for Oncology (v0.4.0)}

The \texttt{digital-twins/} platform implements patient-specific modeling across three subdirectories:
\begin{itemize}
  \item \textbf{Patient Modeling} (\texttt{patient\_model.py}): Six tumor types with four growth kinetics models---exponential ($V(t) = V_0 e^{kt}$), logistic ($V(t) = \frac{K}{1 + (\frac{K}{V_0} - 1)e^{-rt}}$), Gompertz ($V(t) = K \cdot (\frac{V_0}{K})^{e^{-\alpha t}}$), and von Bertalanffy
  \item \textbf{Treatment Simulation} (\texttt{treatment\_simulator.py}): Multi-modality treatment including chemotherapy, radiation, immunotherapy, targeted therapy, and combination protocols with RECIST 1.1 response assessment~\cite{eisenhauer2009recist}
  \item \textbf{Clinical Integration} (\texttt{clinical\_integrator.py}): EHR, PACS, LIS, RIS, and TPS system integration with a federated clinical bridge
\end{itemize}

The visualization infrastructure (v0.9.3) provides 30 Plotly-based interactive charts across three thematic sets for publication-quality rendering of digital twin data.

\subsection{Technology Stack and Dependencies}

The platform is designed for minimal dependency footprint while supporting optional advanced capabilities. Table~\ref{tab:techstack} summarizes the technology stack across required and optional tiers.

\begin{table}[H]
\centering
\caption{Technology stack with version requirements and integration status}
\label{tab:techstack}
\small
\begin{tabular}{@{}llll@{}}
\toprule
\textbf{Package} & \textbf{Version} & \textbf{Purpose} & \textbf{Status} \\
\midrule
\multicolumn{4}{l}{\textit{Required (Python 3.10+, no GPU)}} \\
NumPy & $\geq$1.24.0 & Array computation, federated MLP & Integrated \\
SciPy & $\geq$1.11.0 & ODE solvers, optimization & Integrated \\
scikit-learn & $\geq$1.3.0 & ML algorithms & Integrated \\
pandas & $\geq$2.0.0 & Data manipulation & Integrated \\
cryptography & $\geq$41.0.0 & HMAC-SHA256, encryption & Integrated \\
PyYAML & $\geq$6.0 & Configuration parsing & Integrated \\
pytest & $\geq$7.4.0 & Testing framework & Integrated \\
ruff & ==0.15.1 & Linting \& formatting (pinned) & CI-enforced \\
\midrule
\multicolumn{4}{l}{\textit{Optional: Deep Learning \& Medical Imaging}} \\
PyTorch & $\geq$2.5.0 & Deep learning models & Bridge ready \\
MONAI & $\geq$1.3.0 & Medical image analysis & Optional \\
pydicom & $\geq$2.4.0 & DICOM parsing & Optional \\
\midrule
\multicolumn{4}{l}{\textit{Optional: Agentic AI Frameworks}} \\
LangChain & $\geq$0.3.0 & LLM chains & Interface \\
CrewAI & $\geq$0.80.0 & Multi-agent orchestration & Interface \\
MCP & $\geq$1.1.0 & Model Context Protocol & Interface \\
\midrule
\multicolumn{4}{l}{\textit{Optional: Simulation Frameworks}} \\
NVIDIA Isaac Sim & 4.2.0 & High-fidelity surgical sim & Bridge \\
MuJoCo & $\geq$3.1.0 & Physics simulation & Bridge \\
Gazebo & Harmonic & ROS 2 integration & Guide \\
PyBullet & $\geq$3.2.6 & Lightweight simulation & Guide \\
\bottomrule
\end{tabular}
\end{table}

\subsection{Federated Coordinator Implementation Details}

The core federated learning coordinator (\texttt{federated/coordinator.py}, v0.2.0) implements three aggregation strategies with convergence detection. The following code illustrates the FedProx proximal term, which penalizes client models that deviate too far from the global model---critical for non-IID hospital data distributions:

\begin{lstlisting}[caption={FedProx aggregation with proximal regularization from \texttt{federated/coordinator.py} (v0.2.0)}]
class FederatedCoordinator:
    """Orchestrates federated training rounds with
    FedAvg, FedProx, or SCAFFOLD aggregation."""

    def aggregate_fedprox(self, client_updates, global_weights):
        """FedProx: weighted average with proximal term mu."""
        total_samples = sum(u["n_samples"] for u in client_updates)
        aggregated = {}
        for key in global_weights:
            weighted_sum = sum(
                u["weights"][key] * (u["n_samples"] / total_samples)
                for u in client_updates
            )
            # Proximal term: mu * (w - w_global)
            proximal = self.mu * (weighted_sum - global_weights[key])
            aggregated[key] = weighted_sum - proximal
        return aggregated
\end{lstlisting}

The data harmonization module (\texttt{federated/data\_harmonization.py}, v0.2.0, $\sim$10,161 bytes) provides DICOM/FHIR vocabulary mapping and unit conversion across hospital sites, ensuring that heterogeneous data formats are standardized before federated aggregation. The site enrollment manager (\texttt{federated/site\_enrollment.py}, v0.2.0, $\sim$11,523 bytes) implements quality-weighted site selection for multi-site trial coordination.

\subsection{CLI Tools for Trial Operations (v0.4.0)}

Five command-line tools provide operational support for trial engineers:

\begin{enumerate}
  \item \textbf{DICOM Inspector} (\texttt{tools/dicom-inspector/}): Inspection and validation of DICOM medical images with tag verification and JSON output
  \item \textbf{Dose Calculator} (\texttt{tools/dose-calculator/}): Radiation and chemotherapy dose calculation with NCCN/RTOG reference constants and BSA-based dosing
  \item \textbf{Trial Site Monitor} (\texttt{tools/trial-site-monitor/}): Multi-site monitoring with GREEN/YELLOW/RED status classification and enrollment tracking
  \item \textbf{Sim Job Runner} (\texttt{tools/sim-job-runner/}): Cross-framework simulation job launcher with automatic framework detection
  \item \textbf{Deployment Readiness} (\texttt{tools/deployment-readiness/}): Checklist-based validation scoring with PASS/FAIL/WARNING/REQUIRES\_VERIFICATION status levels
\end{enumerate}

% =================================================================
\section{Discussion}
\label{sec:discussion}
% =================================================================

\subsection{Unified Platforms as a Precondition for Physical AI Adoption}

The PAI Oncology Trial FL platform demonstrates that unified infrastructure---combining federated learning, privacy, regulation, cross-framework interoperability, and multi-organization cooperation---is a \emph{necessary precondition} for transitioning the oncology industry to using robots in physical AI clinical trials. Without such unification:
\begin{itemize}
  \item \textbf{Privacy fragmentation} would require each trial site to independently implement HIPAA Safe Harbor de-identification, increasing the risk of compliance gaps
  \item \textbf{Regulatory inconsistency} would delay FDA submissions as each site uses different tracking systems with incompatible audit trails
  \item \textbf{Simulation incompatibility} would prevent cross-site validation of robotic procedures, as policies trained in Isaac Sim cannot directly execute in MuJoCo or Gazebo without the unification bridge
  \item \textbf{Aggregation vulnerability} would expose federated models to data poisoning without standardized secure aggregation and differential privacy across all participating sites
\end{itemize}

The platform's 235 Python modules provide a single coherent answer to these fragmentation risks, enabling the industry transition from software-only AI to physical AI in oncology trials.

\subsection{Triple AI Peer Review as Quality Assurance Methodology}

The sequential Codex-to-Claude Code peer review pipeline (v0.9.4--v0.9.9) establishes a replicable methodology for AI code quality assurance. The correction trajectory---Cycle 1 (CI/process) $\to$ Cycle 2 (compliance/security) $\to$ Cycle 3 (hardening/secrets)---demonstrates systematic risk elimination where each cycle addresses progressively deeper issues. All six review/fix releases were completed on 2026-02-19, a pace impractical for human reviewers processing 86,800 LOC~\cite{kawchak2026peerreview}.

The dual-manufacturer approach mitigates single-AI bias: Codex (OpenAI) reviewing Claude Code (Anthropic) ensures neither manufacturer's training data blind spots persist in the final codebase. This is analogous to dual-auditor requirements in financial compliance, where independence between auditors strengthens assurance.

\subsection{Agentic AI for Clinical Trial Operations}

The six agentic AI examples (v0.5.0) demonstrate that LLM agents can be safely integrated into clinical trial workflows when constrained by:
\begin{enumerate}
  \item \textbf{Safety gates} (Script 05): IEC 80601-2-77 aligned pre/post-condition checks with human-in-the-loop approval for high-risk actions~\cite{iec80601}
  \item \textbf{Protocol grounding} (Script 06): RAG-based retrieval from regulatory documents with faithfulness verification and citation tracking
  \item \textbf{Audit trails} (Script 01): 21 CFR Part 11 compliant logging of all LLM tool invocations with hash-chain integrity~\cite{cfr_part11}
  \item \textbf{Privacy enforcement} (All scripts): PHI detection and de-identification applied before any data leaves the local site
\end{enumerate}

\subsection{Comparison with Existing Approaches}

Compared to existing federated learning platforms for oncology (Dana-Farber multi-center platform~\cite{danafarberfederated2025}, DigiONE~\cite{digione2024}), PAI Oncology Trial FL is unique in integrating:
\begin{itemize}
  \item Physical AI robotics (surgical robot interfaces, sensor fusion, cross-framework simulation)
  \item Complete regulatory submission pipeline (eCTD compilation, multi-jurisdiction intelligence)
  \item Agentic AI capabilities (MCP, ReAct, RAG for clinical trial operations)
  \item Digital twins with four tumor growth kinetics models
  \item Triple AI peer review for code quality assurance
\end{itemize}

\subsection{Quantitative Development Velocity}

The development velocity achieved by AI-driven development is quantifiable from the release history:

\begin{table}[H]
\centering
\caption{AI development velocity metrics by date}
\label{tab:velocity}
\small
\begin{tabular}{@{}lcrl@{}}
\toprule
\textbf{Date} & \textbf{Releases} & \textbf{Est.\ LOC Added} & \textbf{Key Achievements} \\
\midrule
2026-02-17 & v0.1.0--v0.5.0 & $\sim$55,000 & Core platform + 31 examples \\
2026-02-18 & v0.6.0--v0.9.1 & $\sim$25,000 & Privacy, security, analytics \\
2026-02-19 & v0.9.2--v1.0.1 & $\sim$6,800 & Peer review, visualization, docs \\
2026-02-27 & v1.1.0 & $\sim$500 & Paper, documentation updates \\
\midrule
\textbf{Total} & \textbf{21 releases} & $\sim$\textbf{86,800} & \textbf{4 calendar days of AI work} \\
\bottomrule
\end{tabular}
\end{table}

This represents approximately 21,700 LOC per active development day---a rate enabled by Claude Code Opus 4.6's ability to generate production-quality Python with comprehensive docstrings, type annotations, test coverage, and regulatory compliance disclaimers in each prompt cycle. The PAI Oncology Trial FL platform builds upon the original Physical AI Oncology Trials repository~\cite{kawchak2026physical} (51 Python modules, $\sim$40,526 LOC), extending it with a 4.3$\times$ increase in codebase size and the addition of federated learning, clinical analytics, regulatory submissions, and the triple AI peer review quality process.

% =================================================================
\section{Limitations and Future Work}
\label{sec:limitations}
% =================================================================

\subsection{Current Limitations}

\subsubsection{Claude Code Limitations}
\begin{itemize}
  \item \textbf{Context window constraints:} Large codebases require careful prompt decomposition; the 18-prompt development strategy was necessary because no single prompt could address the full 86,800 LOC platform
  \item \textbf{Deterministic output variability:} Despite \texttt{np.random.seed(42)} seeding, Claude Code's stochastic generation may produce slightly different implementations across identical prompts
  \item \textbf{Domain knowledge boundaries:} Clinical oncology details (e.g., specific drug dosing protocols, exact FDA submission timelines) required external verification by Mr.\ Kawchak
  \item \textbf{Testing limitations:} While Claude Code generated 82 test files, property-based testing coverage (via Hypothesis) and integration test depth could be expanded
\end{itemize}

\subsubsection{Codex Limitations}
\begin{itemize}
  \item \textbf{Review-only role:} Codex was restricted to generating recommendations without directly modifying code, limiting its ability to verify fix correctness
  \item \textbf{Recommendation granularity:} Some recommendations (e.g., ``improve exception handling'') required Claude Code to interpret scope and priority
  \item \textbf{Contextual blind spots:} Codex reviews focused on individual files rather than cross-module architectural patterns, potentially missing systemic issues
\end{itemize}

\subsubsection{Human Oversight Limitations}
\begin{itemize}
  \item \textbf{Prompt engineering effort:} Mr.\ Kawchak authored 18 prompts over the development lifecycle; scaling this approach requires prompt engineering expertise
  \item \textbf{Review bottleneck:} All AI-generated PRs required human review and merge approval, creating a serial bottleneck in the development pipeline
  \item \textbf{Clinical validation gap:} The platform is research-grade only; clinical deployment requires independent IRB approval, FDA clearance, and prospective clinical validation
\end{itemize}

\subsection{Future Work}

Three development proposals are documented in \texttt{DEVELOPMENT\_PROPOSALS.md}~\cite{kawchak2026fl}:
\begin{enumerate}
  \item \textbf{Real-Time Federated Safety Signal Detection:} Pharmacovigilance algorithms (PRR, ROR, BCPNN, MGPS, SPRT), MedDRA classification, ICH E2B reporting
  \item \textbf{Genomic Biomarker Integration Pipeline:} VCF/MAF variant processing, ACMG/AMP classification, TMB/MSI/HRD scoring, federated GWAS with privacy preservation
  \item \textbf{Adaptive Trial Design Engine:} Bayesian inference, response-adaptive randomization, Thompson sampling, platform/master protocol designs
\end{enumerate}

% =================================================================
\section{Conclusion}
\label{sec:conclusion}
% =================================================================

This paper presents the PAI Oncology Trial FL platform (v1.1.0)---a comprehensive federated learning framework for physical AI oncology trials comprising 235 Python modules ($\sim$86,800 LOC) developed across 21 releases over 11 days. The platform unifies five critical infrastructure pillars: privacy (18 HIPAA Safe Harbor identifiers), regulation (FDA/IRB/GCP/multi-jurisdiction), cross-framework interoperability (Isaac Sim/MuJoCo/Gazebo/PyBullet), standards and benchmarking, and multi-organization cooperation.

End-to-end workflow demonstrations across 31 example scripts---including six production-grade agentic AI implementations using MCP, ReAct, real-time monitoring, autonomous orchestration, safety-constrained execution, and RAG-based compliance---establish that the platform supports the full clinical trial lifecycle from federated model training through regulatory submission.

The triple AI peer review process (v0.9.4--v0.9.9) resolved 31/31 code recommendations at 100\% completion through sequential Codex-to-Claude Code review-fix cycles, demonstrating that dual-manufacturer AI code review can operate at speeds and depths impractical for traditional human-only review of codebases exceeding 86,000 LOC. Combined with the 61 security audit findings resolved in v0.7.0, the platform has undergone 92 total quality corrections establishing a verifiable trust chain.

Unified platforms such as PAI Oncology Trial FL are a necessary precondition for transitioning the oncology industry to using robots in physical AI clinical trials. The open-source availability of this platform (MIT license) enables the research community to build upon, validate, and extend these capabilities toward clinical deployment.

\textbf{DOI:} \href{https://doi.org/10.5281/zenodo.18795507}{10.5281/zenodo.18795507}

% =================================================================
\section*{References}
% =================================================================

\bibliographystyle{unsrt}
\bibliography{references}

% =================================================================
\section*{Acknowledgments}
% =================================================================

The author would like to acknowledge Anthropic for providing access to Claude Code, Claude Cowork; and OpenAI for providing access to ChatGPT.

% =================================================================
\section*{Ethical Disclosures}
% =================================================================

The author of the article declares no competing interests.

% =================================================================
\section*{Rights and Permissions}
% =================================================================

This article is distributed under the terms of the Creative Commons Attribution 4.0 International License (CC BY 4.0), which permits unrestricted use, distribution, and reproduction in any medium, provided the original author(s) and source are properly credited, a link to the Creative Commons license is provided, and any modifications made are indicated. To view a copy of this license, visit \url{https://creativecommons.org/licenses/by/4.0/}.

% =================================================================
\section*{Cite This Article}
% =================================================================

Kawchak K. Federated Learning Physical AI Oncology Trials Unification. Zenodo. 2026; \href{https://doi.org/10.5281/zenodo.18795507}{10.5281/zenodo.18795507}.

\end{document}
